% !PRINTED 43
\weberpage{43}
%<<CONT
\[
S_0 = \int \frac{dx}{\sqrt{f(x)}}, \quad S_2 = \int \frac{x^2\,dx}{\sqrt{f(x)}}, \quad S_{-1} = \int \frac{dx}{(x-\alpha)\sqrt{f(x)}},
\]
w\"ahrend $S_1 = \displaystyle\int \frac{x\,dx}{\sqrt{f(x)}}$ durch die Substitution $x^2 = y$ auf das nicht elliptische Integral
\[
\frac{1}{2}\int \frac{dy}{\sqrt{y^2 + a_2 y + a_4}}
\]
zur\"uckgef\"uhrt wird. Hier ist $S_0$ von der ersten, $S_2$ von der zweiten, $S_{-1}$ von der dritten Gattung. F\"ur die \gesperrt{Legendresche} Normalform setzen wir $f(x) = (1-x^2)(1-\varkappa^2 x^2)$ und nehmen als Normalintegral der zweiten Gattung nicht $S_2$ selbst, sondern $S_0 - \varkappa^2 S_2$. Dadurch ergeben sich die Normalintegrale der drei Gattungen:
\[
\int \frac{dx}{\sqrt{(1-x^2)(1-\varkappa^2 x^2)}}, \quad \int \sqrt{\frac{1-\varkappa^2 x^2}{1-x^2}}\,dx, \quad \int \frac{dx}{(x-\alpha)\sqrt{(1-x^2)(1-\varkappa^2 x^2)}}
\]
und wenn man $x = \sin\varphi$ setzt:
\[
\int \frac{d\varphi}{\sqrt{1-\varkappa^2\sin^2\varphi}}, \quad \int \sqrt{1-\varkappa^2\sin^2\varphi}\,d\varphi, \quad \int \frac{d\varphi}{(\sin\varphi - \alpha)\sqrt{1-\varkappa^2\sin^2\varphi}}.
\]

Dasselbe erh\"alt man, wenn man das elliptische Differential in der Normalform
\[
\frac{dz}{\sqrt{z(1-z)(1-\varkappa^2 z)}}
\]
annimmt und nach dem Falle II verf\"ahrt.

\begin{center}
\textbf{\S\ 13. Das \gesperrt{Additionstheorem.}}
\end{center}

Das von \gesperrt{Euler} entdeckte \gesperrt{Additionstheorem} der elliptischen Integrale besteht in dem Satze, der f\"ur die ganze weitere Theorie von fundamentaler Bedeutung ist, da\ss, wenn von den drei Wertpaaren
\[
x_1, \sqrt{f(x_1)}; \quad x_2, \sqrt{f(x_2)}; \quad x_3, \sqrt{f(x_3)}
\]
zwei als gegeben vorausgesetzt werden, man auf \gesperrt{algebraischem Wege} das dritte so bestimmen kann, da\ss\ die Differentialgleichung
\weq{1}{\frac{dx_1}{\sqrt{f(x_1)}} + \frac{dx_2}{\sqrt{f(x_2)}} + \frac{dx_3}{\sqrt{f(x_3)}} = 0}
befriedigt ist. Darin ist $f(x)$ eine gegebene Funktion vom dritten oder vierten Grade. Es ist dies ein spezieller Fall des gro\ss en
%CONT>>
